% critique.tex — f170r, vv. 636–678 : commentaire critique
% Articles textuels tirés du commentaire de l'auteur (latex/Aristophane-Thesmo Harmattan 2023 *.tex),
% texte et notes conservés ; répartition détaillée dans repartition.md.
\cc{637–40}{κἄπειτ᾽ ἀποδύσετ᾽ ἐννέα παίδων μητέρα; ... καὶ νὴ Δία τίτθος γ᾽ ὥσπερ ἡμεῖς οὐκ ἔχει.}{Le manuscrit de Ravenne porte en marge du vers 638 μν\auDessus{η}{σ}, indiquant ainsi que Mnésiloque prend la parole ; le manuscrit ne présente alors plus aucun signe de changement de locuteur jusqu'au vers 641\nc{Il est d'usage dans notre manuscrit que le \italique{signe} soit un παράγραφος semblable à un tiret cadratin ` --- ' lorsque le changement de locuteur se fait en début de vers. Si le changement intervient en milieu de vers, il y a plus souvent un \latin{punctum duplex} ` : ' ou bien une abréviation indiquant le nom du personnage qui prend la parole. Par abus de langage, j'utiliserai le terme de \italique{paragraphos} pour désigner indifféremment l'un de ces signes.}. Or Mnésiloque ne peut aucunement prononcer le vers 639, commentaire sur sa poitrine étonnamment plate (ὡς καὶ στιβαρά τις φαίνεται καὶ καρτερά).
Il y a donc un \italique{paragraphos} à rétablir.
Je dis \italique{un} car les éditeurs suivent quasi-unanimement Brunck en `rétablissant' deux : au début du vers 638 (Clisthène) et au début du vers 639 (la femme -- Mika -- reprend la parole).
Une autre solution me semble plus légère : laisser à Mnésiloque les deux vers 637-638 et à Mika les vers 639-640.
Cette retouche comporte surtout l'avantage d'être bien meilleure quant à l'action et au jeu de scène.\\
Presque tous les éditeurs depuis Brunck, disais-je, retirent le vers 638 à Mnésiloque pour l'accorder à Clisthène. Double erreur ! \ \ \ \italique{Primo}, parce qu'ils sont alors contraints de surajouter une seconde émendation : transformer le χάλᾳ du manuscrit en un χάλα (impératif lancé par Clisthène) alors que le iota adscrit est nettement écrit. \ \ \ \italique{Secundo}, parce que Mika, au vers 636, avait déjà ordonné à Clisthène de déshabiller Mnésiloque (ἀπόδυσον αὐτόν) ; il n'y a aucune raison de douter qu'il ne soit en train de lui obéir. Un peu d'action, que diable ! une pièce d'Aristophane ne doit pas avoir la même pesanteur que l'administration française. (Selon nos éditeurs, Mika ordonnerait à Clisthène de déshabiller Mnésiloque qui transmettrait alors l'ordre à l'intéressé... qui finirait par ne rien faire ?..)\\
Le vers 638, χαλᾷ ταχέως τὸ στρόφιον ὦναίσχυντε σύ, était pourtant excellent dans la bouche de Mnésiloque : Clisthène est en train de le déshabiller, il feint l'indignation (χαλᾷ est un subjonctif dit \italique{délibératif}) devant l'empressement (ταχέως) que cet homme -- car Clisthène reste un homme -- met à enlever le \italique{soutien gorge} d'une femme. « Il défait mon strophion avec promptitude. Impudent ! » Mnésiloque continue de jouer une femme -- vertueuse même, mais un peu tard. Qu'y a-t-il de plus limpide ? l'apostrophe ὦναίσχυντε σύ conclue logiquement la tirade de mère effarouchée qu'affecte le peu crédible Mnésiloque (et qui fait d'ailleurs écho au précédent ἀναίσχυντός τις εἶ (v. 611) lancé par le même Mnésiloque au même Clisthène\nc{Quelques éditeurs avaient même prétexté de ce double ἀναίσχυντος pour attribuer l'ἀναίσχυντός τις εἶ du vers 611 aussi à Clisthène, histoire que leurs multiples tripatouillages présentent une apparence d'homogénéité...}).}
\cc{642–3}{νῦν· τότε δὲ μήτηρ ἦσθα παίδων ἐννέα. / ἀνίστασ᾽ ὀρθός. ποῖ τὸ πέος ὠθεῖς κάτω;}{Il n'y a aucun \italique{paragraphos} aux vers 642 et 643.
Pourtant, ce n'est assurément pas Mnésiloque (qui a la parole au vers précédent) qui peut prononcer ces vers. On ajoute souvent deux changements d'interlocuteurs : Mika pour le vers 642 et Clisthène pour le vers 643. Par simplicité, je préfère n'en ajouter qu'un et attribue donc le distique entier à Clisthène.}
\cc{655–86}{Le chœur compris entre les vv. 655-686.}{Pour ce qui touche à la colométrie, je reprends celle de Lætitia\nc{\italique{The songs of Aristophanes}, pp. 418-422.}. (Le passage des vers 667-686 a de fortes correspondances métriques avec les vers 707-725).}
\cc{656}{τῶν θ᾽ ἱματίων ἀποδύσας.}{Willems, fidèle à son poste de \italique{gardien de la tradition}, justifie contre les retoucheurs de texte l'expression ἀποδῦναι $+$ gén. (C'est sa \nieme{20} note aux \italique{Thesmophories}.)}
\cc{657}{εἰσελήλυθεν.}{La leçon du manuscrit de Ravenne est \italique{contra metrum}. Tannegui Le Fèvre (Faber dans l'apparat), le père de Mme Dacier, avait proposé la simple correction ἐσελήλυθε. Austin\nc{Voir ses notes : « Textual problems in Ar. \italique{Thesm}. » (1987) ; je reprends la trame de celle qui nous intéresse présentement.} la réfuta : ἐς ne précède jamais une voyelle chez Aristophane. Restent les conjectures de Fritzsche et de Handley. Fritzsche proposa de lire ἀνελήλυθε, l'expliquant par une haplographie ἀνὴρ (ἀν)ἐλήλυθε suivie d'une complétion postérieure εἰσ-. Handley proposa ἐπελήλυθε (la confusion entre ΙC et Π en onciales étant une erreur courante) -- hypothèse beaucoup plus simple et, partant, plus satisfaisante.}
\cc{664}{εἴ τις ἐν τόποις ἑδραῖος.}{Qu'avons-nous affaire d'ajouter, avec MM. Blaydes et van Leeuwen, τούτοισι τοῖς devant τόποις ? La même locution adverbiale se lit dans un fragment de Nausicratès : [τρίγκλαι] ἃς κλύδων Αἰξωνικὸς ἐν τόποις παιδεύεται (\italique{Fragm. Com. Gr}., t. IV, p. 575 [= Athénée 325e = Nausicratès, fr. 1.6-11 Kassel-Austin ; le texte n'est pas établi avec certitude : on trouve aussi : ἐντόπους à la place de ἐν τόποις. Même si l'exemple n'est pas probant, Willems n'en a pas moins raison ; Lætitia remarque que l'adjonction de Blaydes accroît la disparité entre les deux stances des vers 663-666 -- là où l'on chercherait à rétablir une symétrie.]) et la locution inverse se lit dans Sophocle : ἥτις μ᾽ ἀποικίσειεν ἐκ τόπων (\italique{Trach}., 955), où jamais personne n'a trouvé à redire (ἐκ τούτων τῶν τόπων.). Je ne vois pas que ces tournures soient plus bizarres que ἐν ποσί et ἐκ ποδῶν, ou bien que ἐν χώρᾳ et ἐκ χώρας, πρὸ ou παρὰ χειρῶν, etc. (Note de Willems.)}
